%!TEX root = ../main.tex
\chapter{Introduction}\label{introduction}

Music has always been a social phenomenon, and that's never been more true than in the internet age, with discussion and discovery being easier than ever before. In the streaming era, the internet has become an important hub for both music listening and discovery, as services have been created not just for people to play music, but to do things like write reviews, catalog their collections, or chronicle their listening histories. 

Larger web services like the sales platform Discogs or the listening tracker Last.fm make data from their platforms available to others, allowing an ecosystem of smaller web apps run by independent developers to proliferate using their data, which is for the most part freely accessible. This project will concern the Last.fm service and its associated web app ecosystem.

\section{About Last.fm}

Last.fm is a web service that collects music listening data through an API. Users connect their Last.fm account to a supported music player like iTunes or Spotify, which sends data about each song played through the software. This data is used to build a comprehensive record of everything they've listened to while connected. The site uses this data to create statistics, visuals, and tools with the data it receives, which will be expanded on throughout this section.

Every song played by a user creates an entry in Last.fm's database. Each contains a song title, artist name, and timestamp, with optional entries for an album name and album artist. They call these entries "scrobbles", after the name given to the software Last.fm runs on, the Audioscrobbler. An example of how these scrobbles are displayed on a user's profile is shown in \autoref{scrobbles}

\begin{figure}[!ht]
	\begin{center}
		\woopic{scrobbles}{.7}
		%\vspace{-.2 in}
		\caption{Four scrobbles on a user's profile, each track displaying their album, title, artist, and timestamp}\label{scrobbles}
	\end{center}
\end{figure}

Entries can be deleted by their user at any time, but can only be added within two weeks of the current data and time. This makes data from more than 14 days ago slightly more trustworthy, as new entries cannot be added after that time.

\subsection{Artist profiles}

While the main purpose of Last.fm is to build profiles of users' listening data, the site also builds profiles of each artist, album, and track that gets submitted. Each profile gets its own dedicated webpage that contains an editable wiki, a six-month history of how many unique listeners the artist/album/track had on each day (shown in \autoref{6month}), and other user-generated widgets like a schedule of live performances for artists and a gallery of cover art for albums. These pages are auto-generated the first time an artist/album/track is submitted by a user and automatically refresh their information daily

\begin{figure}[!ht]
	\begin{center}
		\woopic{6monthgraph}{.9}
		%\vspace{-.2 in}
		\caption{An example of Last.fm's graph of artist's six-month listening history}\label{6month}
	\end{center}
\end{figure}

%[NB: most of this info is gathered from my firsthand experience and some footnotes on the website. Thus, it isn't really cite-able.]

\subsection{Last.fm's "Mainstream Score"}

Last.fm delivers listening reports to its users every week describing their listening trends over that period of time. While all users receive these reports, paying users get access to a few extra stats and charts, including the Mainstream Score, a metric telling users how popular the artists they've been listening to are. The tool uses a largely unspecified method to assign a popularity weight to each artist, and provides the user a normalized score out of 100 and the identity of the most obscure artist they played in that time.

\begin{figure}[!ht]
	\begin{center}
		\woopic{mainstreamscore}{.45}
		%\vspace{-.2 in}
		\caption{An example of how Last.fm represents its Mainstream Score}\label{mainstreamscore}
	\end{center}
\end{figure}

While this feature provides interesting information to users, it has a few key weaknesses. First, it is only available for paying users, those having a Pro subscription to the site. Second, it offers no further information than the two basic points of information shown in \autoref{mainstreamscore} -- even the heights of the bars of the histogram are hard-coded and do not change from week to week. Data about the popularity of each individual artist is being collected by the site, but is not available to users in any way.

The aim of this project is to create a web app that uses scraping -- a process to remotely obtain information from webpages -- to collect publicly available data on the popularity of certain artists and compare it with listening data supplied by a given user. This will improve on the Mainstream Score in two key ways: first, it will be free and open-source, with the source code uploaded to GitHub for anyone to view or modify as they see fit; second, it provides more information to the user than Last.fm provides, ideally showing the user which artists contributed most to the weighting of their score. The site will take in as input the number of times a user has listened to different artists in a chosen timeframe and will output data about how mainstream the music taste represented by the input data is. A goal of the project is to fit into the wider ecosystem of Last.fm web apps while being the first of its kind to integrate input data from both scraping.

This paper will be laid out as follows: \autoref{bg} will cover background information on the relevant topics of web scraping and Web APIs, \autoref{fullstack} discusses the standards and practices of how web apps are designed and deployed, \autoref{databases} discusses the design and structure of databases and how web apps use them for persistent data storage, \autoref{dataviz} discusses the data visualization principles behind the visualizations the project will use and the nature of the tools used in designing charts and graphics for web apps, \autoref{implement} discusses my implementation of the software using the tools and concepts laid out in previous chapters, and \autoref{future} discusses future plans for adding further functionality to the final software. Finally, \autoref{conclusion} concludes the paper.